\makeatletter
\def\input@path{{./}{../}{../../}{../../../}{../../../../}{preamble/}{../preamble/}{../../preamble/}{../../../preamble/}{../../../../book2026-2027/preamble/}}
\makeatother

\input{preamble_exams.tex}

\title{[Activity title]}
\author{}
\date{}

\begin{document}
\maketitle

\setcounter{tocdepth}{1}
\setcounter{secnumdepth}{2}
\phantomsection
\label{toc}
\tableofcontents

\section*{Evaluated criteria}
\addcontentsline{toc}{section}{Evaluated criteria}
\begin{description}
    \item[\textbf{C[n]}] [Criterion description.]
\end{description}

% Optional introduction. Include only the preparation students need.
% \section*{Introduction}
% \addcontentsline{toc}{section}{Introduction}
% [Definitions, notation, assumptions, or reusable method.]
% [General or symbolic worked example, when useful.]

% Repeat this section pattern if grouping problems improves the progression.
\section{Worked Solution Examples}
\label{sec:worked-examples}

\SectionProblemsTableOfContents{
    \SectionProblemLink{prob:example}{[Problem title]}
}

\subsection{[Problem title]}
\label{prob:example}

[Self-contained scenario with information, units, and constraints.]

% Replace [n] with the actual criterion suffix. The command supplies C.
\QuestionsTasks[sec:worked-examples]{
    {[n]}/{sol:example-task}/{[Criterion-linked task prompt.]}
}

\subsubsection*{[Worked solution title]}
\label{sol:example-task}

[Map the scenario to the mathematical model and state the rule or formula.]

[Show the substitutions and intermediate reasoning.]

[Interpret the result in context.]

\secInicio[prob:example][sec:worked-examples]

% Optional independent practice. Use a separate group and place its solutions
% as specified in the brief. Give any linked solutions unique labels.

\end{document}
